\documentclass[11pt]{article}
\usepackage{amsmath,amssymb,amsthm}
\usepackage{geometry}
\usepackage{booktabs}
\usepackage{hyperref}
\usepackage{algorithm}
\usepackage{algpseudocode}
\usepackage{listings}
\usepackage{xcolor}
\usepackage{tcolorbox}
\geometry{margin=1in}

\newtheorem{theorem}{Theorem}
\newtheorem{lemma}{Lemma}
\newtheorem{definition}{Definition}
\newtheorem{corollary}{Corollary}

\lstset{
  basicstyle=\small\ttfamily,
  breaklines=true,
  frame=single,
  columns=fullflexible,
  keywordstyle=\bfseries,
}

\title{Building an Implicit-Dictionary LZSS Compressor \\
\large A Practical Guide to Sliding-Window Matching, Optimal Parsing, \\
and Where the Real Gains Over DEFLATE/zstd/LZMA Come From}
\author{}
\date{}

\begin{document}
\maketitle

\begin{tcolorbox}[colback=yellow!5!white,colframe=yellow!75!black,title=How to use this document]
This is written as a build guide, not just a theory reference. Each section ends with a
concrete deliverable: a data structure, an algorithm, or a file to write. Follow the
sections in order -- each one only depends on what came before it. Section
\ref{sec:calibration} at the end is the most important one to internalize \emph{before}
you start coding: it tells you honestly where the compression-ratio gains over
DEFLATE/zstd/LZMA actually come from, so you don't spend weeks on the part that matters
least.
\end{tcolorbox}

\tableofcontents

\section{Goal and Non-Goal}

\textbf{Goal}: an LZSS-style compressor where the dictionary is never stored. The decoder
reconstructs it implicitly, byte by byte, from the same data it has already decoded --
exactly the same way the encoder does. Nothing about a "dictionary" is ever written to
the compressed file; only literal bytes and (length, distance) back-references are.

\textbf{Non-goal}: reproducing DEFLATE. DEFLATE's specific parameters (32KB window,
greedy/lazy matching, static per-block Huffman) were 1996 engineering trade-offs for
speed and memory, not a template for maximum ratio. This guide deliberately departs from
several of them; each departure is called out explicitly.

\section{Why No Header Dictionary Is Needed: The Causality Argument}
\label{sec:causality}

\begin{definition}[Implicit dictionary]
An implicit dictionary at position $l$ is simply the already-decoded prefix
$x_{1..(l-1)}$ of the output itself. There is no separate structure to store or
transmit.
\end{definition}

\begin{theorem}[Decoder can always reconstruct the dictionary on its own]
\label{thm:causal}
If the encoder only ever emits back-references whose source position $p$ satisfies
$p < l$ (the current position), the decoder can resolve every reference using only
bytes it has \emph{already} produced, in the same left-to-right order the encoder used.
No side information about the dictionary is required.
\end{theorem}

\begin{proof}
By induction on $l$. At $l = 1$ nothing has been decoded, so only a literal can be
emitted (there is nothing to reference). Suppose the decoder has correctly reconstructed
$x_{1..(l-1)}$. If the encoder emits a literal at $l$, the decoder appends it directly.
If the encoder emits a back-reference $(s, \delta)$ with source position
$p = l - \delta$, then $p < l \le l$, so $x_p$ (and, byte by byte, $x_p, \dots,
x_{p+s-1}$, even if $p + s - 1 \ge l$, since each output byte becomes available to be
read the instant it is produced) is already in the decoder's output buffer. The decoder
copies it forward exactly as the encoder did. By induction, this holds for all $l$.
\end{proof}

This is the same argument that dissolved the "per-path hash" confusion from earlier in
this project: match validity is a property of the \emph{fixed input string}, not of
which parse/path was chosen to reach a given position. Section~\ref{sec:matchfinder}
turns this into an algorithm for finding valid matches quickly.

\begin{tcolorbox}[colback=blue!5!white,colframe=blue!50!black,title=What this replaces]
Your current \texttt{compressed\_header.c} / \texttt{code\_map.c} store an explicit
dictionary. Under the design in this document, those become unnecessary: the header
only needs to store per-block metadata (block length, entropy-coder tables if any), not
a symbol table. This is the single biggest structural simplification in this rewrite.
\end{tcolorbox}

\section{Sliding Window: Choosing Its Size}
\label{sec:window}

\begin{center}
\begin{tabular}{@{}lll@{}}
\toprule
Compressor & Window size & Note \\
\midrule
DEFLATE & 32{,}768 bytes & Fixed by RFC 1951's bit-field widths \\
zstd (level $\le$ 19) & up to 8~MB & Configurable \\
zstd \texttt{--ultra -22} & up to 128~MB & \\
LZMA/xz & up to $1.5$~GB & Configurable, dictionary-size parameter \\
\bottomrule
\end{tabular}
\end{center}

\textbf{Recommendation}: do not hardcode a small window. Use the entire block you are
compressing as the window (your current \texttt{BLOCK\_SIZE} of 64{,}000 bytes is already
larger than DEFLATE's whole window). If you later support multi-block or streaming
files, make the window at least as large as \texttt{BLOCK\_SIZE}, ideally the size of the
whole file for offline compression -- there is no ratio reason to shrink it, only a
memory-budget reason, and modern machines make that budget generous.

\section{Match Finding}
\label{sec:matchfinder}

\subsection{Data structure: hash chains}

\begin{definition}[Hash chain matchfinder]
For every position $i$, hash the next $H$ bytes ($H=3$ or $4$ is standard) into a table
\texttt{head[hash]} pointing to the most recent position with that hash, with each
position also storing \texttt{prev[i]} pointing to the next-most-recent earlier position
sharing the same hash -- a singly linked list per hash bucket, walked nearest-first.
\end{definition}

\begin{algorithm}[H]
\caption{Insert position $i$ into the hash-chain structure}
\begin{algorithmic}[1]
\Procedure{InsertPosition}{$i$, $data$, $head[\,]$, $prev[\,]$}
  \State $h \gets \text{Hash}(data[i..i+H-1])$
  \State $prev[i] \gets head[h]$
  \State $head[h] \gets i$
\EndProcedure
\end{algorithmic}
\end{algorithm}

\begin{algorithm}[H]
\caption{Find match candidates ending at position $l$ (equivalently, starting at $i = l -
s + 1$ for the position under consideration)}
\label{alg:findmatches}
\begin{algorithmic}[1]
\Procedure{FindMatches}{$i$, $data$, $head[\,]$, $prev[\,]$, $max\_chain$}
  \State $h \gets \text{Hash}(data[i..i+H-1])$
  \State $p \gets head[h]$
  \State $best\_len \gets 0$
  \State $candidates \gets \emptyset$ \Comment{list of (length, distance), one per achievable length}
  \State $tries \gets 0$
  \While{$p \ge 0$ and $tries < max\_chain$}
    \State $len \gets \text{MatchLength}(data, p, i)$ \Comment{extend while bytes match, capped by window}
    \If{$len > best\_len$}
      \For{$s \gets best\_len+1$ to $len$}
        \State $candidates[s] \gets (s,\, i - p)$ \Comment{nearest distance for every new length reached}
      \EndFor
      \State $best\_len \gets len$
    \EndIf
    \State $p \gets prev[p]$
    \State $tries \gets tries + 1$
  \EndWhile
  \State \textbf{return} $candidates$
\EndProcedure
\end{algorithmic}
\end{algorithm}

\begin{tcolorbox}[colback=blue!5!white,colframe=blue!50!black,title=Key property this relies on]
Because the chain is walked nearest-first, the \emph{first} time a given length is
reached during the walk, its distance is automatically the smallest (cheapest) distance
achieving that length. You never need to compare distances for the same length --
just keep walking until longer matches stop appearing, or the chain/budget runs out.
\end{tcolorbox}

\subsection{How deep to walk the chain}

\texttt{max\_chain} trades encode speed for match quality. DEFLATE's higher settings use
depths in the hundreds; since your goal is ratio, not encode speed, use a generous depth
(thousands, or unbounded for smaller blocks like your current 64{,}000-byte block size --
exhaustive chain walking on 64KB is entirely affordable on modern hardware).

\subsection{Upgrading the matchfinder later: binary trees and suffix automata}

Hash chains can still miss the best match if the chain is very long and gets truncated by
\texttt{max\_chain}. Two standard upgrades, in order of typical benefit-to-effort:

\begin{itemize}
  \item \textbf{Binary-tree matchfinder} (what LZMA's strongest mode, \texttt{BT4}, uses):
  replace each hash bucket's linked list with a binary search tree ordered by the
  matched suffix, giving exact nearest-distance-per-length without a depth cutoff, at
  $O(\log n)$ insertion/search instead of $O(\text{chain length})$.
  \item \textbf{Suffix automaton / suffix array}: gives you the true longest previous
  factor at every position in one linear-time pass over the whole block, as discussed
  earlier in this project. More work to implement correctly, but removes match-finding
  quality as a variable entirely.
\end{itemize}

Start with hash chains (Algorithm~\ref{alg:findmatches}); it is by far the easiest to get
correct, and with a generous \texttt{max\_chain} on a 64KB block it will already be very
close to exhaustive.

\section{Candidate Representation: Retiring the Old Graph Structure}
\label{sec:candidates}

Your existing \texttt{Graph}/\texttt{GraphNode} machinery enumerates every possible
chunk length $1..k$ at every position (the "virtual parent link" design), because the
old dictionary scheme needed to consider every candidate length to know which ones
repeat anywhere in the file. A real matchfinder does not need this: at each position it
hands you a short, explicit list of \emph{real} candidates.

\begin{definition}[Candidate list]
For each position $l = 1, \dots, n$, define $C_l$ as the list of pairs $(s, \delta)$
returned by \textsc{FindMatches} -- typically a handful of entries, not up to $k=128$.
\end{definition}

\textbf{Recommendation}: replace \texttt{Graph}/\texttt{GraphNode}/\texttt{MAX\_GRAPH\_NODES}
with a flat, per-position array of candidates:
\begin{lstlisting}[language=C]
typedef struct {
    uint32_t length;
    uint32_t distance;
} MatchCandidate;

typedef struct {
    MatchCandidate candidates[MAX_CANDIDATES_PER_POSITION];
    uint8_t count;
} PositionCandidates;

PositionCandidates candidates[BLOCK_SIZE]; // one entry per position
\end{lstlisting}
This is smaller, simpler, and removes \texttt{SEQ\_LENGTH\_LIMIT} as an artificial cap on
match length -- a real match can be as long as the matchfinder actually finds (DEFLATE's
258 and LZMA's 273 caps are themselves just wire-format choices, not fundamental
limits; choose a cap based on your entropy coder's length-encoding range, not before).

\section{Cost Model}
\label{sec:cost}

You already have the right building block: \texttt{lzss\_cost.h}'s
\texttt{lzss\_match\_cost\_bits(length, distance)} and
\texttt{lzss\_literal\_cost\_bits()}. Nothing about this section changes once real
distances are available -- the cost model was already written distance-aware, precisely
in anticipation of this rewrite.

\begin{tcolorbox}[colback=blue!5!white,colframe=blue!50!black,title=Reminder: why this is exactly the right unit now]
Under the old dictionary/header scheme, a chunk's true cost depended on how many
distinct dictionary entries the whole file ended up needing -- a shared, global,
history-dependent cost the flat DP could not see, which is why the DP lost to greedy in
your benchmark. Under pure back-references, a match's cost is a function of its own
$(length, distance)$ alone, nothing shared with any other token. This is exactly the
assumption the DP's optimality proof requires, so the same DP that lost to greedy
before will not have that problem here.
\end{tcolorbox}

\section{Optimal Parsing}
\label{sec:parsing}

\subsection{The flat DP (exactly optimal, given the cost model above)}

This is precisely Section~7 of \texttt{doc/writeup.tex}, restated for candidate lists
instead of graph levels:

\begin{algorithm}[H]
\caption{Optimal parse via forward DP}
\begin{algorithmic}[1]
\Procedure{OptimalParse}{$data$, $candidates[\,]$, $n$}
  \State $d[0] \gets 0$
  \For{$l \gets 1$ to $n$}
    \State $d[l] \gets d[l-1] + \textsc{LiteralCostBits}()$
    \State $back[l] \gets \textsc{Literal}$
    \For{each $(s, \delta) \in candidates[l]$}
      \State $cand \gets d[l - s] + \textsc{MatchCostBits}(s, \delta)$
      \If{$cand < d[l]$}
        \State $d[l] \gets cand$
        \State $back[l] \gets (\textsc{Match}, s, \delta)$
      \EndIf
    \EndFor
  \EndFor
  \State \Return \Call{ReconstructPath}{$back$, $n$}
\EndProcedure
\end{algorithmic}
\end{algorithm}

\textbf{Complexity}: $O(n \cdot \bar{c})$ time, $O(n)$ space, where $\bar{c}$ is the
average candidate-list length (typically small, bounded by \texttt{max\_chain} or by how
many distinct lengths were actually found -- not by $k$). This is the direct successor
of \texttt{shortest\_path\_dp.c}: same forward-sweep structure, simpler inner loop (no
\texttt{get\_parent\_level}/\texttt{useless} bookkeeping needed), same optimality proof.

\subsection{Extension: repeat-distance modeling (practical, not exactly optimal)}
\label{sec:repdist}

Real data reuses the \emph{same} back-reference distance repeatedly (column-aligned
data, periodic structure). LZMA exploits this with 4 cached "recent distances"
(\texttt{rep0..rep3}); matching one of them costs far fewer bits than encoding a fresh
distance.

\begin{tcolorbox}[colback=red!5!white,colframe=red!60!black,title=Honesty check]
Tracking recent distances exactly turns the DP state into (position, exact set of
recently used distances) -- and distances range up to the window size, so this state
space is not small in general. \emph{Exact} joint optimization of parse and
repeat-distance cache is not tractable this way, and there is no widely known
polynomial algorithm for it in full generality (it is closely related to the "smallest
grammar" family of hard problems flagged earlier in this project). Real encoders,
including LZMA's own optimal-parse mode, do not solve this exactly either -- they use a
bounded-lookahead approximation. What follows is that same standard, practical
approximation, not a new exact algorithm.
\end{tcolorbox}

\textbf{The practical approximation} (what LZMA and zstd's optimal-parse modes actually
do): extend the DP state to $(l, \texttt{rep0})$, where \texttt{rep0} is \emph{only the
single most recently used distance along the best path reaching $l$} (not all four, and
not all combinations) -- carried forward from \texttt{back[l]} the same way
\texttt{shortest\_path\_dp.c} already carries forward a single winning predecessor per
level. At each position, in addition to the matchfinder's candidates, check whether
\texttt{rep0} (inherited from the predecessor state) still matches at this position; if
so, offer it as an extra candidate with a cheaper cost (no full distance encoding
needed, just a "reuse rep0" flag):

\begin{algorithm}[H]
\caption{Optimal parse with a single repeat-distance slot}
\begin{algorithmic}[1]
\Procedure{OptimalParseWithRep0}{$data$, $candidates[\,]$, $n$}
  \State $d[0] \gets 0$; $rep[0] \gets 0$
  \For{$l \gets 1$ to $n$}
    \State $d[l] \gets d[l-1] + \textsc{LiteralCostBits}()$;\ $back[l] \gets \textsc{Literal}$;\ $rep[l] \gets rep[l-1]$
    \For{each $(s, \delta) \in candidates[l]$}
      \State $cand \gets d[l - s] + \textsc{MatchCostBits}(s, \delta)$
      \If{$cand < d[l]$}
        \State $d[l] \gets cand$;\ $back[l] \gets (\textsc{Match}, s, \delta)$;\ $rep[l] \gets \delta$
      \EndIf
    \EndFor
    \State $r \gets rep[l-1]$
    \If{$r > 0$ and \Call{MatchLengthAt}{$data$, $l - r$, $l$} $> 0$}
      \State $s' \gets$ \Call{MatchLengthAt}{$data$, $l-r$, $l$}
      \State $cand \gets d[l - s'] + \textsc{RepMatchCostBits}(s')$ \Comment{cheaper: no distance field}
      \If{$cand < d[l]$}
        \State $d[l] \gets cand$;\ $back[l] \gets (\textsc{RepMatch}, s')$;\ $rep[l] \gets r$
      \EndIf
    \EndIf
  \EndFor
  \State \Return \Call{ReconstructPath}{$back$, $n$}
\EndProcedure
\end{algorithmic}
\end{algorithm}

This stays $O(n \cdot \bar{c})$ -- one extra $O(1)$-ish check per position -- and
typically captures most of the practical benefit of repeat-distance caching. Extending
to 2--4 tracked slots (closer to LZMA's rep0..rep3) is a straightforward generalization
of the same idea if benchmarking shows it is worth the added complexity; start with one
slot and measure before adding more.

\section{Entropy Coding Integration}
\label{sec:entropy}

Your plan (item 3: range coding / context mixing) is already the right choice --
strictly more expressive than DEFLATE's Huffman stage, which is limited to integer-bit
code lengths per symbol.

\subsection{The cost-model / entropy-coder feedback loop}

The bit costs in Section~\ref{sec:cost} are \emph{estimates} (an Elias-gamma-style
approximation), not what your adaptive entropy coder will actually spend, since its
costs depend on model state that evolves with the actual token sequence chosen. The
standard, well-tested resolution (used by LZMA/zstd) is iterative:

\begin{enumerate}
  \item Parse once using the static cost estimates from Section~\ref{sec:cost}.
  \item Feed the resulting token sequence through the real entropy coder; measure the
  actual bits it assigned to each literal/length/distance symbol.
  \item Re-derive per-symbol cost estimates from those measured statistics.
  \item Re-run \textsc{OptimalParse} (Section~\ref{sec:parsing}) with the updated costs.
  \item Repeat 2--4 rounds, or until total size stops improving.
\end{enumerate}

This is not a proof of global optimality (exact joint optimization of parse and adaptive
model state is intractable in general, for the same reason as
Section~\ref{sec:repdist}) -- but it is an empirically excellent fixed point, and it is
exactly what today's best general-purpose compressors ship.

\subsection{Context choices for the entropy stage, roughly in order of expected payoff}

\begin{enumerate}
  \item Separate probability models for literals vs.\ match-length vs.\ distance
  symbols (do not share one model across all three).
  \item Condition literal probabilities on the preceding 1--2 bytes (order-1/order-2
  context), not just an unconditional byte frequency table.
  \item Condition on whether the previous token was a literal or a match (DEFLATE and
  LZMA both find this genuinely useful).
  \item If you pursue full context mixing: combine several such predictors with a
  logistic mixer (PAQ/cmix-style) rather than picking one context order -- this is where
  most of the ratio gap to state-of-the-art compressors actually lives, more so than in
  parsing sophistication.
\end{enumerate}

\section{Implementation Roadmap}
\label{sec:roadmap}

\begin{center}
\begin{tabular}{@{}p{0.5cm}p{4.3cm}p{7.5cm}@{}}
\toprule
\# & Deliverable & Notes \\
\midrule
1 & Hash-chain matchfinder & Algorithm~\ref{alg:findmatches}. New file, e.g.\
\texttt{matchfinder.c}. \\
2 & Candidate-list data structure & Section~\ref{sec:candidates}. Replaces
\texttt{Graph}/\texttt{GraphNode} for this purpose. \\
3 & Flat optimal-parse DP & Adapt \texttt{shortest\_path\_dp.c} to iterate over
candidate lists; reuse \texttt{lzss\_cost.h} as-is (already written for this). \\
4 & Round-trip test harness & Verify byte-exact decompression on your existing
\texttt{tests/} corpus before optimizing further. \\
5 & Benchmark vs.\ gzip/zstd/xz & Establish your real baseline gap \emph{before}
adding complexity -- see Section~\ref{sec:validation}. \\
6 & Repeat-distance (rep0) DP extension & Section~\ref{sec:repdist}. Only after 1--5
are working and measured. \\
7 & Entropy coder + refinement loop & Section~\ref{sec:entropy}. Likely your single
biggest remaining ratio lever -- see Section~\ref{sec:calibration}. \\
\bottomrule
\end{tabular}
\end{center}

\section{Validation and Benchmarking Protocol}
\label{sec:validation}

\begin{enumerate}
  \item \textbf{Correctness first, always}: every change gets a full round-trip
  (compress $\to$ decompress $\to$ byte-diff) over your entire \texttt{tests/} corpus,
  including \texttt{tests/cantrbry}, before you look at compression ratio at all.
  \item \textbf{Compare against real baselines on the same files}: \texttt{gzip -9},
  \texttt{zstd --ultra -22}, and \texttt{xz -9e} on every test file, not just your
  previous version -- your own history-vs-history comparisons (DP vs.\ greedy) are
  useful for internal regressions, but the external baselines are what tells you whether
  you are actually competitive.
  \item \textbf{Ablate one change at a time}: matchfinder quality, DP vs.\ greedy,
  rep-distance, entropy stage -- toggle each independently so you know which one is
  actually earning its keep, the same way this document's earlier greedy-vs-DP
  benchmark isolated exactly one variable.
  \item \textbf{Track encode time too}: an exhaustive \texttt{max\_chain} and a full
  context-mixing entropy stage can get slow; note the trade-off even if ratio is your
  primary goal, so you know what you're spending to get it.
\end{enumerate}

\section{Calibration: Where the Real Gains Actually Come From}
\label{sec:calibration}

Read this section before writing code, not after.

\begin{enumerate}
  \item \textbf{Matchfinder quality} (Section~\ref{sec:matchfinder}): a shallow hash
  chain vs.\ an exhaustive one, or vs.\ a binary tree, changes what matches you can even
  see. Real, but bounded, gains.
  \item \textbf{Optimal vs.\ greedy parsing} (Section~\ref{sec:parsing}): you already
  measured this empirically on your own data once the cost model is fair (see the
  earlier greedy-vs-DP benchmark in this project) -- meaningful, but this alone will
  bring you roughly to where LZMA/zstd's own optimal-parse modes already are, not past
  them, since they do this too.
  \item \textbf{Repeat-distance modeling} (Section~\ref{sec:repdist}): often the single
  best ratio-per-effort win on structured/periodic data, and a real point of leverage
  over compressors that skip it.
  \item \textbf{Entropy coding sophistication} (Section~\ref{sec:entropy}): once parsing
  is near-optimal, this is realistically where most of the remaining gap to
  state-of-the-art ratios lives -- how well you predict the next symbol's probability,
  not how you choose chunk boundaries. The very best general-purpose compressors (cmix,
  paq-family) get their ratio almost entirely from rich, mixed context models here, not
  from unusual match-finding.
\end{enumerate}

\begin{tcolorbox}[colback=yellow!5!white,colframe=yellow!75!black,title=Honest expectation setting]
Beating gzip/DEFLATE outright is a realistic, achievable target with just
Sections~\ref{sec:matchfinder}--\ref{sec:parsing} (better matchfinder + optimal parse
already exceeds DEFLATE's greedy/lazy design). Matching or narrowly beating zstd at high
levels and competing with LZMA/xz is achievable but will require Sections
\ref{sec:repdist}--\ref{sec:entropy} done well, since those tools already run
near-optimal parsers themselves -- your edge over them has to come from matching or
exceeding their entropy-coding sophistication, not from parsing alone. Decisively
beating LZMA/zstd/xz \emph{in general}, across all data types, is a very high bar few
tools clear; beating them on a \emph{specific} data type you care about (by exploiting
structure they don't model) is a much more tractable and still valuable goal. Aim
there first.
\end{tcolorbox}

\end{document}
